\begin{proof}
\textbf{(1) Degenerate and edge cases.}
\begin{enumerate}
\item[(a)] \textit{All weight on a single index.}  
If $\lambda_k=1$ and $\lambda_i=0$ for $i\neq k$, then
\[
f\!\Bigl(\sum_{i=1}^{n}\lambda_i x_i\Bigr)=f(x_k)=\sum_{i=1}^{n}\lambda_i f(x_i),
\]
so equality holds.

\item[(b)] \textit{Some weights equal zero.}  
If $\lambda_j=0$ for some $j$, the term $\lambda_jx_j$ disappears from the weighted sum and $\lambda_jf(x_j)$ disappears from the right–hand side.  Removing those terms yields an instance of \eqref{jensen} with fewer indices, so the inequality is unchanged.

\item[(c)] \textit{The case $n=1$.}  
When there is only one point, \eqref{jensen} reads $f(x_1)\le f(x_1)$, which is true with equality.
\end{enumerate}
Hence the theorem holds in all degenerate situations.  In what follows we assume $n\ge2$ and that at least two weights are positive.

\medskip

\textbf{(2) Base case $n=2$.}  
Let $\lambda_1=t$, $\lambda_2=1-t$ with $t\in[0,1]$.  The weighted average is $tx_1+(1-t)x_2$, and by convexity \eqref{convex}
\[
f\bigl(tx_1+(1-t)x_2\bigr)\le t f(x_1)+(1-t)f(x_2),
\]
which is exactly \eqref{jensen} for $n=2$.  Thus Jensen’s inequality holds for two points.

\medskip

\textbf{(3) Inductive step.}  
Assume that \eqref{jensen} holds for any collection of $k\ge2$ points with non‑negative weights summing to one.  We prove it for $n=k+1$.

Let $\lambda_1,\dots,\lambda_{k+1}$ satisfy \eqref{weights} and set
\[
\mu:=\sum_{i=1}^{k}\lambda_i\qquad(0\le\mu\le1).
\tag{4}
\]
If $\mu=0$, then $\lambda_{k+1}=1$ and we are in the one‑weight case already treated; hence we may assume $\mu>0$.

Define the combined point
\[
y:=\frac{1}{\mu}\sum_{i=1}^{k}\lambda_i x_i .
\tag{5}
\]
Since the coefficients $\lambda_i/\mu$ are non‑negative and sum to one, $y$ is a convex combination of the first $k$ points.

\medskip

\textit{3.1 Two‑point convexity applied to $(y,x_{k+1})$.}  
Using \eqref{4}–\eqref{5} we rewrite the full weighted average as
\[
\sum_{i=1}^{k+1}\lambda_i x_i
= \mu y + (1-\mu)x_{k+1}.
\tag{6}
\]
Applying convexity \eqref{convex} with $t=\mu$, $x=y$, $y=x_{k+1}$ gives
\[
f\!\bigl(\mu y+(1-\mu)x_{k+1}\bigr)\le \mu f(y)+(1-\mu)f(x_{k+1}).
\tag{7}
\]

\medskip

\textit{3.2 Bounding $f(y)$ by the induction hypothesis.}  
Since $y$ is a convex combination of $x_1,\dots,x_k$, the induction hypothesis (with normalized weights $\lambda_i/\mu$) yields
\[
f(y)=
f\!\Bigl(\sum_{i=1}^{k}\frac{\lambda_i}{\mu}x_i\Bigr)
\le
\sum_{i=1}^{k}\frac{\lambda_i}{\mu}\,f(x_i).
\tag{8}
\]
Multiplying \eqref{8} by the positive $\mu$ we obtain
\[
\mu f(y)\le\sum_{i=1}^{k}\lambda_i f(x_i).
\tag{9}
\]

\medskip

\textit{3.3 Completion of the inductive step.}  
Insert \eqref{9} into \eqref{7} and use \eqref{6} for the argument of $f$:
\[
f\!\Bigl(\sum_{i=1}^{k+1}\lambda_i x_i\Bigr)
\le
\sum_{i=1}^{k}\lambda_i f(x_i)+(1-\mu)f(x_{k+1}).
\tag{10}
\]
Because of \eqref{4}, $1-\mu=\lambda_{k+1}$, and \eqref{10} becomes
\[
f\!\Bigl(\sum_{i=1}^{k+1}\lambda_i x_i\Bigr)
\le
\sum_{i=1}^{k+1}\lambda_i f(x_i),
\]
which is precisely \eqref{jensen} for $n=k+1$.

Thus the inequality for $k$ points implies the inequality for $k+1$ points.

\medskip

\textbf{(4) Conclusion by induction.}  
The base case $n=2$ has been proved, and the inductive step shows that validity for $k$ points implies validity for $k+1$ points.  By the principle of mathematical induction, Jensen’s inequality
\[
f\!\Bigl(\sum_{i=1}^{n}\lambda_i x_i\Bigr)\le\sum_{i=1}^{n}\lambda_i f(x_i)
\tag{3}
\]
holds for every integer $n\ge1$.

\medskip

\textbf{(5) Equality conditions.}
\begin{itemize}
\item \textit{General convex $f$.}  
Equality in \eqref{convex} for a pair $(x,y)$ and $t\in(0,1)$ occurs iff $f$ is affine on the interval $[x,y]$.  In the inductive proof, equality in \eqref{7} forces $f$ to be affine on the segment joining $y$ and $x_{k+1}$, while equality in \eqref{8} forces $f$ to be affine on $\operatorname{conv}\{x_1,\dots,x_k\}$.  Consequently, equality in \eqref{jensen} holds \emph{iff} $f$ is affine on the convex hull of $\{x_i:\lambda_i>0\}$.

\item \textit{Strictly convex $f$.}  
If $f$ is strictly convex, the only way an affine restriction can occur on a non‑trivial interval is when that interval collapses to a point.  Hence equality in \eqref{jensen} can occur only when all points with positive weight coincide, i.e.\ there exists $c\in\mathbb R$ such that $x_i=c$ whenever $\lambda_i>0$.
\end{itemize}
In particular, equality also occurs when at most one weight $\lambda_i$ is non‑zero (the inequality reduces to an identity).

\medskip

Therefore, Jensen’s inequality holds for any convex function and any finite collection of real numbers with non‑negative weights summing to one, with the equality cases described above.
\end{proof}